
\section{Justificativa}


A abordagem do problema de prever o comportamento de compra dos consumidores em e-commerce se justifica pela crescente importância de compreender as dinâmicas de consumo no contexto digital. O mercado online é altamente competitivo e dinâmico, e empresas que conseguem antecipar as intenções de compra de seus clientes têm uma vantagem significativa. A capacidade de prever quando um cliente realizará uma nova compra ou quando pode abandonar a plataforma permite que as empresas ajustem suas estratégias de marketing de forma mais eficiente, direcionando campanhas personalizadas, otimizando o estoque e melhorando a experiência do cliente. Além disso, prever a rotatividade de clientes (churn) é crucial para reduzir os custos com aquisição de novos consumidores e aumentar a retenção dos clientes existentes, que são fontes mais lucrativas para os negócios.

A utilização da análise de sobrevivência é particularmente adequada para esse tipo de estudo, pois permite modelar o tempo até a ocorrência de um evento específico, como a realização de uma compra ou o abandono de um cliente. Modelos como o de Cox são eficazes para identificar os fatores que influenciam o comportamento do consumidor ao longo do tempo, permitindo a inclusão de variáveis como histórico de compras, frequência de visitas e características demográficas. Além disso, essa abordagem leva em consideração a censura, ou seja, a situação em que o evento de interesse (compra ou abandono) não ocorre dentro do período de observação.

A abordagem do tempo de vida via algoritmos de árvore de sobrevivência, ou árvore de decisão para dados de sobrevivência e floresta aleatória também para este mesmo tipo de dados é uma abordagem possível e promissora principalmente atualmente quando os bancos de dados em geral são grandes.

Dessa forma, a análise de sobrevivência não apenas oferece uma metodologia robusta para a previsão de comportamentos de compra, mas também permite uma compreensão mais aprofundada dos fatores que influenciam a lealdade e a rotatividade dos clientes. \cite{sobrevivencia}